\begin{table}[!htbp]\centering
\caption{Decomposing the dose response: eligibility versus intensive margin}
\label{tab:decomposition}
\begin{adjustbox}{max width=\textwidth}
\begin{threeparttable}
\begin{tabular}{lcc}
\toprule
 & Floor within reach & Offered depth, given reach \\
 & (eligibility margin) & (intensive margin) \\
\midrule
Essential $\times$ compensation price (\$100/MWh) & $-0.086^{***}$ & $0.005$ \\
 & (0.025) & (0.019) \\
Wild-cluster-bootstrap $p$ (Rademacher) & 0.003 & 0.789 \\
Range across June-2022 treatments & $-0.088$ to $-0.076$ & $0.005$ to $0.036$ \\
Randomization-inference $p$ (999 month-permutations) & 0.046 & --- \\
Month-grain WLS slope (19 months, HC1 $p$) & $-0.096$ (0.0012) & --- \\
Observations (base sample) & 140,259 & 28,555 \\
\bottomrule
\end{tabular}
\begin{tablenotes}\footnotesize
\item The eligibility margin is an indicator for the unit's minimum-stable quantum (its frozen floor: 40~MW per Torrens unit) being offered within dispatch's reach; it is identical under both cheap thresholds because the floor block, when offered, sits at the $-\$1{,}000$ band. The intensive margin is the cheap-capacity share on the subsample with the floor in reach (a decomposition aid: it conditions on an outcome). Specification, matching, controls, clustering, and June-2022 treatments exactly as in Table~\ref{tab:rq2}. The cost-indexed intensive margin behaves identically (base $
-0.019
$, $p=
0.55
$). Interpretations for all cells were committed in the project record before estimation.
\item $^{*}p<0.10$, $^{**}p<0.05$, $^{***}p<0.01$ (analytic).
\end{tablenotes}
\end{threeparttable}
\end{adjustbox}
\end{table}
